\documentclass[10pt,a4paper]{amsart}
\usepackage{amsmath,amssymb,amsthm,graphicx}
\usepackage{verbatim}
\newenvironment{reference}{\comment}{\endcomment}
%\newenvironment{reference}{}{}
\newenvironment{slogan}{\comment}{\endcomment}
\newenvironment{history}{\comment}{\endcomment}

% For commutative diagrams we use Xy-pic
\usepackage[all]{xy}

% We use 2cell for 2-commutative diagrams.
\xyoption{2cell}
% Theorem environments.
%
\theoremstyle{plain}
\newtheorem{theorem}[subsection]{Theorem}
\newtheorem{proposition}[subsection]{Proposition}
\newtheorem{lemma}[subsection]{Lemma}

\theoremstyle{definition}
\newtheorem{definition}[subsection]{Definition}
\newtheorem{example}[subsection]{Example}
\newtheorem{exercise}[subsection]{Exercise}
\newtheorem{situation}[subsection]{Situation}

\theoremstyle{remark}
\newtheorem{remark}[subsection]{Remark}
\newtheorem{remarks}[subsection]{Remarks}

\numberwithin{equation}{subsection}

% Macros
%
\def\lim{\mathop{\mathrm{lim}}\nolimits}
\def\colim{\mathop{\mathrm{colim}}\nolimits}
\def\Spec{\mathop{\mathrm{Spec}}}
\def\Hom{\mathop{\mathrm{Hom}}\nolimits}
\def\Ext{\mathop{\mathrm{Ext}}\nolimits}
\def\SheafHom{\mathop{\mathcal{H}\!\mathit{om}}\nolimits}
\def\SheafExt{\mathop{\mathcal{E}\!\mathit{xt}}\nolimits}
\def\Sch{\mathit{Sch}}
\def\Mor{\mathop{\mathrm{Mor}}\nolimits}
\def\Ob{\mathop{\mathrm{Ob}}\nolimits}
\def\Sh{\mathop{\mathit{Sh}}\nolimits}
\def\NL{\mathop{N\!L}\nolimits}
\def\CH{\mathop{\mathrm{CH}}\nolimits}
\def\proetale{{pro\text{-}\acute{e}tale}}
\def\etale{{\acute{e}tale}}
\def\QCoh{\mathit{QCoh}}
\def\Ker{\mathop{\mathrm{Ker}}}
\def\Im{\mathop{\mathrm{Im}}}
\def\Coker{\mathop{\mathrm{Coker}}}
\def\Coim{\mathop{\mathrm{Coim}}}

% Boxtimes
%
\DeclareMathSymbol{\boxtimes}{\mathbin}{AMSa}{"02}

%
% Macros for moduli stacks/spaces
%
\def\QCohstack{\mathcal{QC}\!\mathit{oh}}
\def\Cohstack{\mathcal{C}\!\mathit{oh}}
\def\Spacesstack{\mathcal{S}\!\mathit{paces}}
\def\Quotfunctor{\mathrm{Quot}}
\def\Hilbfunctor{\mathrm{Hilb}}
\def\Curvesstack{\mathcal{C}\!\mathit{urves}}
\def\Polarizedstack{\mathcal{P}\!\mathit{olarized}}
\def\Complexesstack{\mathcal{C}\!\mathit{omplexes}}
% \Pic is the operator that assigns to X its picard group, usage \Pic(X)
% \Picardstack_{X/B} denotes the Picard stack of X over B
% \Picardfunctor_{X/B} denotes the Picard functor of X over B
\def\Pic{\mathop{\mathrm{Pic}}\nolimits}
\def\Picardstack{\mathcal{P}\!\mathit{ic}}
\def\Picardfunctor{\mathrm{Pic}}
\def\Deformationcategory{\mathcal{D}\!\mathit{ef}}
\usepackage{cleveref}
\crefname{originalreference}{original reference}{original references}
\Crefname{originalreference}{Original reference}{Original references}
\makeatletter
\newlabel{lemma-characteristic-vs-minimal-polynomial}{{0BIG}{source}{}{sourcepointer-1}{}}
\newlabel{lemma-characteristic-vs-minimal-polynomial@cref}{{[originalreference][0][]0BIG}{[1][0][]source}{}{}{}}
\makeatother
\crefname{originalreference}{original reference}{original references}
\Crefname{originalreference}{Original reference}{Original references}
\makeatletter
\newlabel{lemma-finite-purely-inseparable}{{09HI}{source}{}{sourcepointer-2}{}}
\newlabel{lemma-finite-purely-inseparable@cref}{{[originalreference][0][]09HI}{[1][0][]source}{}{}{}}
\makeatother
\crefname{originalreference}{original reference}{original references}
\Crefname{originalreference}{Original reference}{Original references}
\makeatletter
\newlabel{lemma-recognize-separable}{{09H3}{source}{}{sourcepointer-3}{}}
\newlabel{lemma-recognize-separable@cref}{{[originalreference][0][]09H3}{[1][0][]source}{}{}{}}
\makeatother
\crefname{originalreference}{original reference}{original references}
\Crefname{originalreference}{Original reference}{Original references}
\makeatletter
\newlabel{lemma-separable-first}{{030K}{source}{}{sourcepointer-4}{}}
\newlabel{lemma-separable-first@cref}{{[originalreference][0][]030K}{[1][0][]source}{}{}{}}
\makeatother
\crefname{originalreference}{original reference}{original references}
\Crefname{originalreference}{Original reference}{Original references}
\makeatletter
\newlabel{lemma-sums-of-powers}{{0EM9}{source}{}{sourcepointer-5}{}}
\newlabel{lemma-sums-of-powers@cref}{{[originalreference][0][]0EM9}{[1][0][]source}{}{}{}}
\makeatother
\crefname{originalreference}{original reference}{original references}
\Crefname{originalreference}{Original reference}{Original references}
\makeatletter
\newlabel{lemma-trace-and-norm-from-minimal-polynomial}{{0BIH}{source}{}{sourcepointer-6}{}}
\newlabel{lemma-trace-and-norm-from-minimal-polynomial@cref}{{[originalreference][0][]0BIH}{[1][0][]source}{}{}{}}
\makeatother
\crefname{originalreference}{original reference}{original references}
\Crefname{originalreference}{Original reference}{Original references}
\makeatletter
\newlabel{lemma-trace-and-norm-tower}{{0BIJ}{source}{}{sourcepointer-7}{}}
\newlabel{lemma-trace-and-norm-tower@cref}{{[originalreference][0][]0BIJ}{[1][0][]source}{}{}{}}
\makeatother
\title{Separability of a finite extension is equivalent to nonvanishing of the trace and nondegeneracy of the trace pairing}
\author{The Stacks Project Authors}
\date{stacks-project:fields.tex:local-snapshot-20261012; GNU FDL 1.2 or later}
\begin{document}
\maketitle
\noindent\textit{Editorial note.} All delivered mathematical text is the author's own, copied from the pinned Stacks Project chapter file fields.tex; no statement or proof step was rewritten, reordered or removed.\par\medskip
\noindent\textit{Editorial note.} The counted proof is the author's complete argument: The trace pairing nondegeneracy statement and the discriminant nonvanishing discussion that follows in the source (source lines 2802 sqq and 2831-2836) build on this lemma but are not part of the counted unit.\par\medskip
\noindent\textit{Editorial note.} The label lemma-separable-trace-pairing names the counted statement, and the proof environment immediately following it in the source is delivered verbatim as its proof.\par\medskip
\noindent\textit{Editorial note.} The counted statement is carried by the author's own lemma environment in the chapter source.\par\medskip
\noindent\textit{Allowed standard prerequisites (editorial).}\begin{itemize}
\item Same chapter: definition-trace-pairing (the symmetric bilinear form Q\_\{L/K\}(alpha, beta) = Trace\_\{L/K\}(alpha beta), source lines 2730-2738), lemma-trace-and-norm-tower (traces compose in towers), lemma-trace-and-norm-from-minimal-polynomial (trace of a power from the minimal polynomial), lemma-separable-first and lemma-finite-purely-inseparable (the purely inseparable degree p layer), lemma-characteristic-vs-minimal-polynomial (minimal = characteristic polynomial for the multiplication map), lemma-recognize-separable (separable iff distinct roots), lemma-sums-of-powers (power sums detect nonvanishing)
\end{itemize}\medskip
\ifdefined\hypertarget\hypertarget{sourcepointer-1}{}\fi
\noindent\textit{Source reference.} Original reference 0BIG: sources/stacks-fields/pinned-original-source.tex, source lines 2616--2616. Prerequisite pointer invoked at 2792; it names a standard fact of the same chapter that the counted argument uses, and every step of the counted proof built on it is retained verbatim in the delivered body.\par\medskip
\ifdefined\hypertarget\hypertarget{sourcepointer-2}{}\fi
\noindent\textit{Source reference.} Original reference 09HI: sources/stacks-fields/pinned-original-source.tex, source lines 1886--1886. Prerequisite pointer invoked at 2775; it names a standard fact of the same chapter that the counted argument uses, and every step of the counted proof built on it is retained verbatim in the delivered body.\par\medskip
\ifdefined\hypertarget\hypertarget{sourcepointer-3}{}\fi
\noindent\textit{Source reference.} Original reference 09H3: sources/stacks-fields/pinned-original-source.tex, source lines 1424--1424. Prerequisite pointer invoked at 2793; it names a standard fact of the same chapter that the counted argument uses, and every step of the counted proof built on it is retained verbatim in the delivered body.\par\medskip
\ifdefined\hypertarget\hypertarget{sourcepointer-4}{}\fi
\noindent\textit{Source reference.} Original reference 030K: sources/stacks-fields/pinned-original-source.tex, source lines 1921--1921. Prerequisite pointer invoked at 2775; it names a standard fact of the same chapter that the counted argument uses, and every step of the counted proof built on it is retained verbatim in the delivered body.\par\medskip
\ifdefined\hypertarget\hypertarget{sourcepointer-5}{}\fi
\noindent\textit{Source reference.} Original reference 0EM9: sources/stacks-fields/pinned-original-source.tex, source lines 1717--1717. Prerequisite pointer invoked at 2799; it names a standard fact of the same chapter that the counted argument uses, and every step of the counted proof built on it is retained verbatim in the delivered body.\par\medskip
\ifdefined\hypertarget\hypertarget{sourcepointer-6}{}\fi
\noindent\textit{Source reference.} Original reference 0BIH: sources/stacks-fields/pinned-original-source.tex, source lines 2643--2643. Prerequisite pointer invoked at 2767; it names a standard fact of the same chapter that the counted argument uses, and every step of the counted proof built on it is retained verbatim in the delivered body.\par\medskip
\ifdefined\hypertarget\hypertarget{sourcepointer-7}{}\fi
\noindent\textit{Source reference.} Original reference 0BIJ: sources/stacks-fields/pinned-original-source.tex, source lines 2713--2713. Prerequisite pointer invoked at 2776; it names a standard fact of the same chapter that the counted argument uses, and every step of the counted proof built on it is retained verbatim in the delivered body.\par\medskip
\begin{lemma}
\label{lemma-separable-trace-pairing}
Let $L/K$ be a finite extension of fields. The following are equivalent:
\begin{enumerate}
\item $L/K$ is separable,
\item $\text{Trace}_{L/K}$ is not identically zero, and
\item the trace pairing $Q_{L/K}$ is nondegenerate.
\end{enumerate}
\end{lemma}

\begin{proof}
It is clear that (3) implies (2). If (2) holds, then pick $\gamma \in L$
with $\text{Trace}_{L/K}(\gamma) \not = 0$. Then if $\alpha \in L$
is nonzero, we see that $Q_{L/K}(\alpha, \gamma/\alpha) \not = 0$.
Hence $Q_{L/K}$ is nondegenerate. This proves the equivalence of
(2) and (3).

\medskip\noindent
Suppose that $K$ has characteristic $p$ and $L = K(\alpha)$ with
$\alpha \not \in K$ and $\alpha^p \in K$. Then $\text{Trace}_{L/K}(1) = p = 0$.
For $i = 1, \ldots, p - 1$ we see that $x^p - \alpha^{pi}$ is the minimal
polynomial for $\alpha^i$ over $K$ and we find
$\text{Trace}_{L/K}(\alpha^i) = 0$ by
Lemma \ref{lemma-trace-and-norm-from-minimal-polynomial}.
Hence for this kind of purely inseparable degree $p$ extension
we see that $\text{Trace}_{L/K}$ is identically zero.

\medskip\noindent
Assume that $L/K$ is not separable. Then there exists a subfield
$L/K'/K$ such that $L/K'$ is a purely inseparable degree $p$ extension
as in the previous paragraph, see
Lemmas \ref{lemma-separable-first} and \ref{lemma-finite-purely-inseparable}.
Hence by Lemma \ref{lemma-trace-and-norm-tower}
we see that $\text{Trace}_{L/K}$ is identically zero.

\medskip\noindent
Assume on the other hand that $L/K$ is separable.
By induction on the degree we will show that
$\text{Trace}_{L/K}$ is not identically zero.
Thus by Lemma \ref{lemma-trace-and-norm-tower} we may assume that
$L/K$ is generated by a single element $\alpha$ (use that if
the trace is nonzero then it is surjective).
We have to show that $\text{Trace}_{L/K}(\alpha^e)$
is nonzero for some $e \geq 0$.
Let $P = x^d + a_1 x^{d - 1} + \ldots + a_d$ be the minimal
polynomial of $\alpha$ over $K$.
Then $P$ is also the characteristic polynomial of the linear
maps $\alpha : L \to L$, see
Lemma \ref{lemma-characteristic-vs-minimal-polynomial}.
Since $L/k$ is separable we see from Lemma \ref{lemma-recognize-separable}
that $P$ has $d$ pairwise distinct roots $\alpha_1, \ldots, \alpha_d$
in an algebraic closure $\overline{K}$ of $K$. Thus these are the
eigenvalues of $\alpha : L \to L$.
By linear algebra, the trace of $\alpha^e$ is
equal to $\alpha_1^e + \ldots + \alpha_d^e$.
Thus we conclude by Lemma \ref{lemma-sums-of-powers}.
\end{proof}

\end{document}
